% Part: first-order-logic
% Chapter: beyond
% Section: intuitionistic-logic

\documentclass[../../../include/open-logic-section]{subfiles}

\begin{document}

\olfileid{fol}{byd}{il}

\olsection{अंतःप्रज्ञावादी तर्कशास्त्र}

द्वितीय-क्रम आणि उच्च-क्रम तर्कशास्त्राच्या उलट, अंतःप्रज्ञावादी प्रथम-क्रम
तर्कशास्त्र हे अभिजात आवृत्तीचे एक निर्बंधित रूप आहे; अधिक “रचनाशील” प्रकारच्या
तर्कविचाराचे प्रतिमान देणे हा त्याचा उद्देश आहे. त्यामागील काही प्रेरणा पुढील
उदाहरणांनी स्पष्ट होतील.

समजा, एखाद्या दिवशी कोणीतरी तुमच्याकडे येऊन सांगितले की त्यांनी अशी नैसर्गिक
संख्या~$x$ ठरवली आहे की $x$ अविभाज्य असेल तर रीमान गृहीतक सत्य असते आणि $x$
संयुक्त असेल तर रीमान गृहीतक असत्य असते. फारच चांगली बातमी!{} रीमान गृहीतक
सत्य आहे की नाही हा गणितातील मोठ्या अनिर्णित प्रश्नांपैकी एक आहे; आणि येथे तर
ही समस्या केवळ गणनेच्या समस्येत, म्हणजे एखादी ठरावीक संख्या अविभाज्य आहे की
नाही हे ठरवण्यात, रूपांतरित झाल्यासारखी दिसते.

$x$ चे ते अद्भुत मूल्य कोणते? ते त्याचे पुढीलप्रमाणे वर्णन करतात: रीमान गृहीतक
सत्य असेल तर $x$ ही $7$ इतकी आणि अन्यथा $9$ इतकी नैसर्गिक संख्या आहे.

तुम्ही चिडून आपले पैसे परत मागता. अभिजात दृष्टिकोनातून वरील वर्णन $x$ चे
एकमेव मूल्य खरोखर ठरवते; पण तुम्हाला प्रत्यक्षात हवे आहे ते
\emph{स्पष्टपणे} दिलेले $x$ चे मूल्य.

आणखी एक, कदाचित कमी कृत्रिम, उदाहरण पाहू. अपरिमेय संख्येचा परिमेय घात घेऊन
परिमेय उत्तर मिळू शकते, हे आपल्याला माहीत आहे. उदाहरणार्थ,
$\sqrt{2}^2 = 2$. अपरिमेय संख्येचा \emph{अपरिमेय} घात घेऊन परिमेय उत्तर
मिळू शकते की नाही, हे तितके स्पष्ट नाही. पुढील प्रमेय याचे होकारार्थी उत्तर
देते:

\begin{thm}
$a^b$ परिमेय असेल अशा $a$ आणि $b$ या अपरिमेय संख्या आहेत.
\end{thm}

\begin{proof}
$\sqrt{2}^{\sqrt{2}}$ चा विचार करा. ही संख्या परिमेय असेल, तर आपले काम झाले:
$a = b = \sqrt{2}$ घेता येईल. अन्यथा ती अपरिमेय आहे. मग
\[
(\sqrt{2}^{\sqrt{2}})^{\sqrt{2}} = \sqrt{2}^{\sqrt{2} \cdot
  \sqrt{2}} = \sqrt{2}^2 = 2,
\]
आणि ही संख्या निश्चितच परिमेय आहे. म्हणून या प्रकरणात $a$ म्हणून
$\sqrt{2}^{\sqrt{2}}$ आणि $b$ म्हणून~$\sqrt 2$ घ्या.
\end{proof}

याला वैध सिद्धता म्हणता येईल का? बहुतेक गणितज्ञांच्या मते होय. पण यात पुन्हा
काहीसे असमाधानकारक असे काही आहे: विशिष्ट गुणधर्म असलेल्या वास्तव संख्यांच्या
जोडीचे अस्तित्व आपण सिद्ध केले, पण ती संख्यांची \emph{कोणती} जोडी आहे हे
सांगता आलेले नाही. हाच निष्कर्ष अशा रीतीने सिद्ध करता येतो की सिद्धतेत
$a$, $b$ ही जोडी \emph{प्रत्यक्ष दिलेली} असते: $a = \sqrt{3}$ आणि
$b = \log_3 4$ घ्या. मग
\[
a^b = \sqrt{3}^{\log_3 4} = 3^{1/2 \cdot \log_3 4} = (3^{\log_3
  4})^{1/2} = 4^{1/2}= 2,
\]
कारण $3^{\log_3 x} = x$.

पहिल्या सिद्धतेतील पाऊल अमान्य असणाऱ्या तर्कविचाराचे प्रतिमान देण्यासाठी
अंतःप्रज्ञावादी तर्कशास्त्र रचले आहे. $!A(x)$ ची पूर्ति करणारा $x$ अस्तित्वात
आहे हे सिद्ध करणे म्हणजे दुसऱ्या सिद्धतेप्रमाणे एखादा ठरावीक~$x$ आणि तो
$!A$ ची पूर्ति करतो याची सिद्धता द्यावी लागते. $!A$ किंवा $!B$ सत्य आहे हे
सिद्ध करण्यासाठी त्यांपैकी एकाची सिद्धता देता आली पाहिजे.

आकारिक भाषेत सांगायचे, तर प्रथम-क्रम तर्कशास्त्राची !!a{derivation} प्रणाली
ठरावीक रीतीने निर्बंधित केल्यावर अंतःप्रज्ञावादी प्रथम-क्रम तर्कशास्त्र मिळते.
त्याचप्रमाणे द्वितीय-क्रम किंवा उच्च-क्रम तर्कशास्त्राच्याही अंतःप्रज्ञावादी
आवृत्त्या असतात. गणितीय दृष्टिकोनातून या केवळ आकारिक निगमन-पद्धती आहेत; पण
आधी नमूद केल्याप्रमाणे, एका प्रकारच्या गणितीय तर्कविचाराचे प्रतिमान देणे हा
त्यांचा उद्देश आहे. हा तर्कविचार गणिताविषयीच्या एखाद्या तात्त्विक दृष्टिकोनाने
(उदाहरणार्थ, ब्रौवर यांच्या अंतःप्रज्ञावादाने) समर्थित मानता येतो; किंवा
बिशप यांच्या रचनाशीलतावादाच्या धर्तीवर अधिक “मूर्त” आणि समाधानकारक असा
गणितीय तर्कविचार मानता येतो; तसेच आकारिक वर्णन अनौपचारिक प्रेरणा खरोखर
पकडते की नाही यावर वादही घालता येतो. पण आपली तात्त्विक भूमिका कोणतीही असो,
अंतःप्रज्ञावादी तर्कशास्त्राचा आकारिक रीतीने मांडलेले तर्कशास्त्र म्हणून
अभ्यास करता येतो; आणि कारणे कोणतीही असोत, अनेक गणितीय तर्कशास्त्रज्ञांना तो
अभ्यास रंजक वाटतो.

अंतःप्रज्ञावादी संयोजकांचे एक अनौपचारिक रचनाशील अर्थनिर्धारण आहे. ते सहसा
BHK अर्थनिर्धारण म्हणून ओळखले जाते (ब्रौवर, हीटिंग आणि कोल्मोगोरोव यांच्या
नावांवरून). ते असे आहे: $!A \land !B$ ची सिद्धता म्हणजे $!A$ ची सिद्धता
आणि $!B$ ची सिद्धता यांची जोडी; $!A \lor !B$ ची सिद्धता म्हणजे $!A$ ची
सिद्धता किंवा $!B$ ची सिद्धता, तसेच यांपैकी नेमके कोणते प्रकरण आहे याची स्पष्ट
माहिती; $!A \lif !B$ ची सिद्धता म्हणजे $!A$ च्या सिद्धतेचे $!B$ च्या
सिद्धतेत रूपांतर करणारी प्रक्रिया; $\lforall[x][!A(x)]$ ची सिद्धता म्हणजे
$x$ च्या कोणत्याही मूल्यासाठी $!A(x)$ ची सिद्धता देणारी प्रक्रिया; आणि
$\lexists[x][!A(x)]$ ची सिद्धता म्हणजे $x$ चे एखादे मूल्य आणि ते मूल्य $!A$
ची पूर्ति करते याची सिद्धता. या अर्थनिर्धारणाचे संगणकीय वर्णन
“करी--हॉवर्ड समरूपण” किंवा “!!{formula}s-हे-प्रकार प्रतिमान” या नावांनी करता
येते: !!a{formula} एखादा विशिष्ट दत्तांश-प्रकार ठरवते असे समजा आणि
सिद्धतांना त्या दत्तांश-प्रकारांच्या संगणकीय वस्तू समजा; या वस्तूंमुळे संबंधित
!!{formula} सत्य आहे हे आपल्याला दिसते.

अंतःप्रज्ञावादी तर्कशास्त्राकडे अनेकदा “विमध्य सिद्धांत वजा केल्यावर उरणारे”
अभिजात तर्कशास्त्र म्हणून पाहिले जाते. पुढील प्रमेय हे अधिक नेमके करते.
\begin{thm}
अंतःप्रज्ञावादी तर्कशास्त्रात पुढील स्वयंसिद्धक-रूपबंध समतुल्य आहेत:
\begin{enumerate}
\item $(\lnot !A \lif \lfalse) \lif !A$.
\item $!A \lor \lnot !A$
\item $\lnot \lnot !A \lif !A$
\end{enumerate}
\end{thm}
कोणत्याही एका रूपबंधापासून इतर दोन्हींचे दृष्टांत मिळवणे हा अंतःप्रज्ञावादी
तर्कशास्त्रातील चांगला सराव आहे.

ब्रौवर यांच्या अंतःप्रज्ञावादाचे आकारीकरण म्हणून मांडलेल्या अंतःप्रज्ञावादी
विधानीय तर्कशास्त्राच्या पहिल्या निगमन-पद्धती कोल्मोगोरोव, ग्लिवेन्को आणि
हीटिंग यांनी स्वतंत्रपणे दिल्या. अंतःप्रज्ञावादी प्रथम-क्रम तर्कशास्त्राचे
(आणि अंतःप्रज्ञावादी गणिताच्या काही भागांचे) पहिले आकारीकरण हीटिंग यांनी
केले. अभिजात तर्कशास्त्रात वैध असलेले अनेक रूपबंध अंतःप्रज्ञावादी
तर्कशास्त्रात वैध नसले, तरी त्यांपैकी अनेक वैधही असतात.

\emph{द्विनकरण रूपांतरण} अभिजात आणि अंतःप्रज्ञावादी तर्कशास्त्रांतील एक
महत्त्वाचा संबंध वर्णिते. ते पुढीलप्रमाणे विगमनाने परिभाषित केले जाते
($!A^N$ हे अभिजात !!{formula}~$!A$ चे “अंतःप्रज्ञावादी” रूपांतर समजा):
\begin{align*}
!A^N & \ident \lnot\lnot !A \quad \text{अण्वीय !!{formula}s $!A$ साठी} \\
(!A \land !B)^N & \ident (!A^N \land !B^N) \\
(!A \lor !B)^N & \ident  \lnot\lnot (!A^N \lor !B^N) \\
(!A \lif !B)^N & \ident (!A^N \lif !B^N) \\
(\lforall[x][!A])^N & \ident \lforall[x][!A^N] \\
(\lexists[x][!A])^N & \ident \lnot\lnot\lexists[x][!A^N]
\end{align*}
कोल्मोगोरोव आणि ग्लिवेन्को यांनी विधानीय तर्कशास्त्रासाठी या रूपांतरणाच्या
आवृत्त्या दिल्या होत्या; विधेय तर्कशास्त्रासाठी ते गोडेल आणि गेंटझेन यांनी
स्वतंत्रपणे दिले. आपल्याला पुढील निष्कर्ष मिळतो.

\begin{thm}
\begin{enumerate}
\item $!A \liff !A^N$ हे अभिजात तर्कशास्त्रात सिद्धतायोग्य आहे.
\item $!A$ हे अभिजात तर्कशास्त्रात सिद्धतायोग्य असेल, तर $!A^N$ हे
  अंतःप्रज्ञावादी तर्कशास्त्रात सिद्धतायोग्य आहे.
\end{enumerate}
\end{thm}

आता पुढील संवादाची कल्पना करता येते. अभिजात गणितज्ञ: “मी $!A$ सिद्ध केले
आहे!” अंतःप्रज्ञावादी गणितज्ञ: “तुमची सिद्धता वैध नाही. तुम्ही प्रत्यक्षात
$!A^N$ सिद्ध केले आहे.” अभिजात गणितज्ञ: “मला तेही चालेल!” अभिजात
गणितज्ञाच्या दृष्टीने अंतःप्रज्ञावादी गणितज्ञ क्षुल्लक भेद करत आहे, कारण ती
दोन्ही समतुल्य आहेत. पण अंतःप्रज्ञावादी गणितज्ञाच्या मते त्यांत फरक आहे.

लक्षात घ्या की वरील रूपांतरण फक्त शुद्ध तर्कशास्त्राशी संबंधित आहे; अभिजात
आणि अंतःप्रज्ञावादी गणितासाठी योग्य \emph{अतार्किक} स्वयंसिद्धके कोणती किंवा
त्यांच्यात कोणता संबंध आहे, या प्रश्नाला ते हात घालत नाही. पण वरील प्रमेयाचा
पुढील किंचित विस्तार काही उपयुक्त माहिती देतो:

\begin{thm}
$\Gamma$ ने $!A$ अभिजात रीतीने सिद्ध होत असेल, तर $\Gamma^N$ ने $!A^N$
अंतःप्रज्ञावादी रीतीने सिद्ध होते.
\end{thm}

दुसऱ्या शब्दांत, काही गृहीतकांपासून $!A$ अभिजात रीतीने सिद्धतायोग्य असेल,
तर त्यांच्या द्विनकरण रूपांतरांपासून $!A^N$ अंतःप्रज्ञावादी रीतीने
सिद्धतायोग्य असते.

एखादे वाक्य किंवा विधानीय !!{formula} अंतःप्रज्ञावादी रीतीने वैध आहे हे
दाखवण्यासाठी फक्त त्याची सिद्धता द्यावी लागते. पण ते वैध नाही हे कसे
दाखवणार? त्यासाठी आपल्याला निर्दोष, आणि शक्यतो संपूर्ण, चिन्हार्थमीमांसा हवी.
क्रिप्के यांनी दिलेली चिन्हार्थमीमांसा हे काम उत्तम रीतीने करते.

अभिजात तर्कशास्त्रासाठी केलेला खेळच येथेही खेळता येतो: चिन्हार्थमीमांसा
परिभाषित करायची आणि निर्दोषता व संपूर्णता सिद्ध करायची. मात्र पुढील भेद लक्षात
घेण्यासारखा आहे. अभिजात तर्कशास्त्रात चिन्हार्थमीमांसा एका अर्थाने
संयोजकांच्या अर्थातच अंतर्भूत असलेली “उघड” चिन्हार्थमीमांसा होती. क्रिप्के
चिन्हार्थमीमांसेसाठी काही अंतःप्रज्ञ स्पष्टीकरण देता येत असले, तरी तिच्यात
तशी अपरिहार्यता जाणवत नाही. शिवाय, अभिजात !!{structure} ही स्वाभाविक गणितीय
संकल्पना आहे; म्हणून !!a{structure} ही संकल्पना अभिजात प्रथम-क्रम
तर्कशास्त्राच्या अभ्यासाचे साधन मानता येते, किंवा अभिजात प्रथम-क्रम
तर्कशास्त्र हे गणितीय !!{structure}s च्या अभ्यासाचे साधन मानता येते.
याउलट, क्रिप्के !!{structure}s कडे फक्त तार्किक रचना म्हणूनच पाहता येते;
त्यांना स्वतंत्र गणितीय महत्त्व आहे असे दिसत नाही.

विधानीय भाषेसाठीची क्रिप्के !!{structure}~$\mModel M = \tuple{W, R, V}$
मध्ये $W$ हा संच, $W$ वरील लघुतम !!{element} असलेला $R$ हा अंशतः क्रम आणि
विधानीय चलांचे $W$ च्या !!{element}s ना केलेले “एकस्वनिक” मूल्यांकन असते.
येथे $W$ चे !!{element}s “जगे” किंवा “ज्ञानावस्था” दर्शवतात; $v \geq u$ हा
घटक $u$ ची “संभाव्य भावी अवस्था” दर्शवतो; आणि $u$ ला दिलेली विधानीय चरे ही
$u$ अवस्थेत सत्य असल्याचे माहीत असलेली विधाने असतात. बलन-संबंध
$\mSat{M}{!A}[w]$ मग हा संबंध भाषेतील कोणत्याही !!{formula}s पर्यंत
विस्तारित करतो; $\mSat{M}{!A}[w]$ चा अर्थ “$w$ अवस्थेत $!A$ सत्य आहे” असा
वाचा. हा संबंध पुढीलप्रमाणे विगमनाने परिभाषित केला जातो:
\begin{enumerate}
\item $\mSat{M}{\Obj p_i}[w]$ तेव्हा आणि केवळ तेव्हाच, जेव्हा $\Obj p_i$
  हे $w$ ला दिलेल्या विधानीय चरांपैकी एक असते.
\item $\mSat/{M}{\lfalse}[w]$.
\item $\mSat{M}{(!A \land !B)}[w]$ तेव्हा आणि केवळ तेव्हाच, जेव्हा
  $\mSat{M}{!A}[w]$ आणि $\mSat{M}{!B}[w]$.
\item $\mSat{M}{(!A \lor !B)}[w]$ तेव्हा आणि केवळ तेव्हाच, जेव्हा
  $\mSat{M}{!A}[w]$ किंवा $\mSat{M}{!B}[w]$.
\item $\mSat{M}{(!A \lif !B)}[w]$ तेव्हा आणि केवळ तेव्हाच, जेव्हा
  $w' \geq w$ आणि $\mSat{M}{!A}[w']$ असताना नेहमी $\mSat{M}{!B}[w']$.
\end{enumerate}
प्रतिउदाहरण देणारी क्रिप्के !!{structure} रचून $\lnot (p \land q) \lif
(\lnot p \lor \lnot q)$ हे अंतःप्रज्ञावादी रीतीने वैध नाही हे दाखवण्याचा
प्रयत्न करणे हा चांगला सराव आहे.

\end{document}
